\cotefile{29}
\scanpage{21}
% archive sheet number: 20
\deleted{$(a_i)_{i\quad I}$ de $\underline{O}_S$, un} une famille $(D_i)_{i\quad I}$ de diviseurs (de Cartier) sur $S$. \unsure{Bour|Pour|\deleted{P}Bour|\deleted{B}Pour} tout $S'$ étale sur $S$, on désigne par $\underline{C}_{S'}$ la catégorie suivante: les objets sont les \deleted{couples triples} \inserted{quadruples $M=$} $(\underline{a},\underline{n},\inserted{G,}\underline{u})$, où $\underline{a}=(a_i)_{i\quad I}$ est une famille d'équations pour les $D_{iS'}$, $\unsure{\underline{n}=\deleted{-}(n_i)_{i\quad I}|\underline{n}=\deleted{(}-n_i)_{i\quad I}}$ une famille d'entiers $1$ premiers aux caractéristiques résiduelles de $S'$, \deleted{et $u:G\quad{}_{\underline{n}}$} $G$ un groupe fini, et $u:G_{S'}\quad{}_{\underline{n}S'}$ un isomorphisme de schémas en groupes sur $S'$. 2tant donné un deuxième tel quadruplet \inserted{$M'=$} $(\underline{a}',\underline{n}',G',u')$, \deleted{un morphisme du second dans le premier est $M'$ dans $M$} $\operatorname{Hom}_{S'}(M',M)$ est vide s\unsure{\deleted{i}a|\deleted{a}i}uf si $\underline{n}'\quad\underline{n}$, auquel cas c'est l'ensemble des $\underline{b}=(b_i)_{i\quad I}$ tels que $\underline{b}^{n}=\underline{a}'\underline{a}^{-1}$. La composition des morphis\unsure{m|n}es,et la définition des foncteurs changements de base pour une flêche $S''\quad S'$ dans $\underline{\mathrm{Et}}_S$, sont clairs. Le foncteur car\unsure{\deleted{é}t|\deleted{t}é}\unsure{\deleted{t}é|\deleted{é}t}sien $\underline{C}\quad(\mathrm{Sch})_{/S}$ est défini en associant à $M=(\underline{a},\underline{n},G,\underline{u})$ le revêtement Kummérien élémentaire $Z_{\underline{a}}^{\underline{n}}$, \deleted{et à une} \deleted{flê} considéré comme un revêtement à groupe d'opérateurs $G$, grâce à l'isomorphisme $G\quad{}_{\underline{n}}$. On définit de faç\deleted{\unsure{a|o|x}}on évidente la valeur de $\underline{r}$ sur les \unsure{\deleted{c}f|\deleted{f}c}lêches de $\underline{C}_{S'}$, et sa compatibilité avec les changements de base. On voit alors qu'on est sous les conditions de l'exemple 3. De plus, le foncteur $\underline{r}$ est fidèle, et il est pleinement fidèle lorsque $S$ est \deleted{noethérien aux points de} localement noethérien, et de plus \unsure{\deleted{j}n|\deleted{i}n|\deleted{n}j|\deleted{n}i}ormal en les points des $\operatorname{supp}D_i$, et si de plus les $D_i$ sont réduits, et toute \inserted{comp\unsure{o|c}sante irréductible d\deleted{e}'un $D_i$ est distincte de toute comp. d'un $D_j$ pour $j\ne i$}

b) Supposons donné un diviseur $D$ sur $S$, \deleted{une famille couvrante $\unsure{S\quad S|S'\quad S}$} et considérons le faisceau \unsure{\deleted{é}s|\deleted{s}é|ś}ur $\underline{\mathrm{Et}}_S$ associé au préfaisceau qui, à tout $S'$, associe l'ensemble des \deleted{diviseurs} ensembles \deleted{$\unsure{\underline{Q}|\underline{\Omega}|\underline{\delta}|\underline{D}|\underline{O}}$} \inserted{$(D_i)$} $\theta$ de diviseurs sur $S'$ ayant les propriétés suivantes: $\unsure{D_{S\overset{=}{'}}=\quad D_i|D_{S\deleted{=}'}=\quad D_i|D_{S\deleted{'}=} =\quad D_i|D_{S\bar{'}}=\quad D_i}$, et pour tout $s'\quad S'$, les restrictions \inserted{\deleted{non nulles}} à $\text{Sêc}(\underline{O}_{S',s'})$ des $D_i$ qui sont non nulles, sont deux à deux disjointes. Soit donnée une section du
